% ==============================================================================
% ПАРТИЯ B03-028: DISS-005-B1564 .. DISS-005-B1573
% Глава 3: Деперсонализация и структура файла конфигурации COMTRADE
% Замечания: DISS-005-C0117 (1 шт.)
% Объекты: 0 шт.
% ==============================================================================

% DISS-005-B1564
\section{Предобработка осциллограмм: деперсонализация и унификация}\label{sec:ch3_preprocessing_depersonalization}

% DISS-005-B1565
После сбора и первичной систематизации осциллограмм следующим этапом является их предобработка. Этот этап направлен на приведение разнородных данных к единому виду, пригодному для дальнейшего анализа и использования, а также на обеспечение конфиденциальности информации.\par\vspace{0.2cm}

% DISS-005-B1566
\subsection{Ключевые аспекты формата COMTRADE для дальнейшей обработки}\label{sec:ch3_comtrade_aspects}

% DISS-005-B1567
Формат COMTRADE (Common Format for Transient Data Exchange for Power Systems) является стандартом IEEE (C37.111) для обмена данными о переходных процессах в энергосистемах. Осциллограммы в этом формате обычно состоят из двух файлов: конфигурационного файла (.cfg) и файла данных (.dat).\par\vspace{0.2cm}

% DISS-005-B1568
\subsubsection{Конфигурационный файл (.cfg)}\label{sec:ch3_comtrade_cfg}

% DISS-005-B1569
Первая строка содержит информацию об объекте, номер регистратора и версию формата. В процессе деперсонализации, описанном далее, информация об объекте и регистраторе из этой строки удаляется. Однако версия формата имеет значение: стандарт COMTRADE, начиная с ревизии 1999 года, предусматривает наличие во второй части файла .cfg (после описания дискретных каналов) информации о коэффициентах трансформации для аналоговых сигналов.\par\vspace{0.2cm}

% DISS-005-B1570
Вторая строка указывает общее количество каналов, а также раздельно количество аналоговых (A) и дискретных (S) каналов. Например, «10,6A,4S». Важно отметить, что в некоторых осциллограммах из нашего датасета общее количество каналов могло не совпадать с суммой аналоговых и дискретных, что потребовало дополнительной проверки при парсинге.\par\vspace{0.2cm}

% DISS-005-B1571
Строки описания аналоговых каналов. Каждая такая строка содержит параметры для одного аналогового канала, разделенные запятыми. Ключевые для нормализации параметры представлены в таблице~1.\par\vspace{0.2cm}

% DISS-005-B1572
Информация о частоте дискретизации и частоте сети располагаются после описания всех каналов и также важны нам для анализа и нормализации.\par\vspace{0.2cm}

% DISS-005-B1573
\paragraph{Таблица\commentnote{DISS-005-C0117}{Алексей Евдаков [2]}{Привести к нормальной стилистики оформления.}~1. Описание параметров аналоговых сигналов.}
\label{tab:ch03_analog_signals_parameters}

% ==============================================================================
% ПАРТИЯ B03-029: DISS-005-B1574 .. DISS-005-B1615
% Глава 3: Таблица 1 (Параметры аналоговых сигналов, TAB-076, 42 ячейки)
% ==============================================================================
% Файл-заполнитель для последующей сборки Таблицы 1 в партии B03-029.

